\section{Ventajas del modelo de flujo y capacidad de generalización}

Tras establecer la relación de equivalencia, el ponente se dirigió a discutir las \textbf{ventajas únicas} del modelo de flujo frente al modelo de difusión —es decir, en qué aspectos el modelo de flujo realmente ``supera'' al modelo de difusión.

\subsection{Mapa de flujo condicional no lineal}

Los modelos de difusión utilizan esencialmente caminos condicionales de forma gaussiana lineal. En cambio, el marco Flow Matching permite definir \textbf{cualquier} mapa de flujo condicional, sin limitarse a la forma lineal:

\begin{itemize}
    \item Se pueden probar mapas no lineales $\psi_t(\mathbf{x}_0 \mid \mathbf{x}_1)$, por ejemplo añadiendo transformaciones no lineales como coseno o sigmoide
    \item Los experimentos demuestran que, aunque los mapas de flujo no lineales actualmente no superan significativamente la forma lineal, la flexibilidad inherente del marco ofrece espacio para futuras exploraciones
\end{itemize}

\subsection{Distribución de referencia arbitraria}

Esta es una de las generalizaciones más importantes del modelo de flujo. Los modelos de difusión requieren que la distribución de referencia sea necesariamente una distribución gaussiana estándar $\mathcal{N}(\mathbf{0}, \mathbf{I})$, ya que toda la teoría (proceso de difusión hacia adelante) se construye sobre la base del ruido gaussiano. En contraste, el modelo de flujo puede utilizar \textbf{cualquier} distribución de referencia:

\begin{importantbox}{Mapeo de distribución a distribución}
Flow Matching permite aprender un mapeo desde \textbf{cualquier} distribución fuente hacia \textbf{cualquier} distribución objetivo. Esto incluye no solo:
\begin{itemize}
    \item Gaussiana estándar $\to$ distribución de datos (tarea de generación clásica)
\end{itemize}
sino también:
\begin{itemize}
    \item Distribución de datos A $\to$ distribución de datos B (transferencia de estilo, conversión de dominios)
    \item Distribución de baja resolución $\to$ distribución de alta resolución (superresolución)
    \item Conversiones entre distribuciones arbitrariamente complejas
\end{itemize}
Esta flexibilidad es algo que el marco del modelo de difusión no posee.
\end{importantbox}

\subsection{Origen ODE vs.\ SDE}

\begin{knowledgebox}{Dos filosofías de modelado}
Otra diferencia fundamental entre los modelos de difusión y los de flujo radica en el \textbf{origen} del modelado:
\begin{itemize}
    \item \textbf{Modelos de difusión}: parten de una \textbf{ecuación diferencial estocástica} (SDE), definiendo el proceso de difusión hacia adelante como un proceso estocástico que añade ruido gaussiano gradualmente; luego se descubre que una ODE \textbf{determinista} equivalente también proporciona la misma distribución marginal
    \item \textbf{Modelos de flujo}: parten directamente de una \textbf{ecuación diferencial ordinaria} (ODE), definiendo un mapeo determinista desde la distribución fuente hacia la distribución objetivo, sin necesidad de introducir aleatoriedad
\end{itemize}
El origen del modelo de flujo es más sencillo y evita las derivaciones complejas de conversión de SDE a ODE.
\end{knowledgebox}

\subsection{Resumen de la sección}

El modelo de flujo generaliza al modelo de difusión en los siguientes aspectos: (1) el mapa de flujo condicional no se limita a formas lineales; (2) la distribución de referencia no se limita a una gaussiana estándar; (3) parte directamente de una ODE, ofreciendo un concepto más sencillo. Estas generalizaciones proporcionan una base teórica para el diseño de modelos generadores más flexibles.
